% ===========================================================================
%  dodatekN_erg_renormalizacja.tex
%  Teoria Generowanej Przestrzeni (TGP) — wersja 1
%
%  DODATEK N: Analiza ERG (Wetterich) dla substratu TGP
%  ------------------------------------------------------
%  Cel:   Uzasadnić postać K(φ)∝φ⁴ jako punkt stały przepływu RG
%         i wyznaczyć λ_eff z równania Wettericha.
%
%  Metoda: Exact Renormalization Group (ERG) w aproksymacji LPA
%          i LPA' (z dynamiczną funkcją kinetyczną).
%
%  Wyniki: K(φ)∝φ⁴ jest zgodny z przepływem ERG przy T~T_c;
%          K_IR/K_UV = 1.13; θ_u = −8.68 (sygnał asymptotycznej
%          bezpieczności); λ_eff wyznaczony numerycznie.
%
%  Implementacja numeryczna: scripts/tgp_erg_wetterich.py
% ===========================================================================

\section*{Dodatek N: Przepływ ERG dla substratu TGP}
\addcontentsline{toc}{section}{Dodatek N: Przepływ ERG}
\label{app:ERG}

\subsection*{N.1 Równanie Wettericha — schemat ogólny}

Dokładne równanie renormalizacyjne (ERG) Wettericha opisuje ewolucję
efektywnego działania $\Gamma_k[\phi]$ ze skalą $k$:
\begin{equation}
  \partial_t \Gamma_k = \frac{1}{2}\,
    \mathrm{Tr}\!\left[
      \bigl(\Gamma_k^{(2)} + R_k\bigr)^{-1}\partial_t R_k
    \right],
  \quad t = \ln(k/\Lambda),
  \label{eq:Wetterich}
\end{equation}
gdzie $\Gamma_k^{(2)}$ to druga pochodna funkcjonalna względem $\phi$,
$R_k(q)$ to regulator infra-czerwony (IR), a $\Lambda$ to skala UV.

\paragraph{Regulator Litmima (optymalny).}
Stosujemy regulator Litmima (Litim 2001):
\begin{equation}
  R_k(q) = (k^2 - q^2)\,\theta(k^2 - q^2),
  \label{eq:Litim_regulator}
\end{equation}
który maksymalizuje konwergencję LPA i pozwala całkować analitycznie
po pędach.

\subsection*{N.2 Aproksymacja lokalna potencjału (LPA)}

W LPA przyjmujemy $\Gamma_k[\phi] = \int [(\nabla\phi)^2/2 + U_k(\phi)]\,d^dx$
(bez renormalizacji funkcji falowej). Dla $d = 3$, $n = 1$
(jedno-składnikowe pole, symetria $\mathbb{Z}_2$) i regulatora Litmima:
\begin{equation}
  \partial_t U_k(\rho) =
    \frac{k^5}{6\pi^2}
    \cdot\frac{1}{k^2 + U_k'(\rho) + 2\rho\,U_k''(\rho)},
  \quad \rho = \tfrac{1}{2}\phi^2.
  \label{eq:LPA_flow}
\end{equation}

\paragraph{Bezwymiarowe zmienne.}
Wprowadzamy zmienne unormowane do skali $k$:
\begin{equation}
  \tilde\rho = \rho\,k^{-(d-2)},
  \quad \tilde u(\tilde\rho) = U_k(\rho)\,k^{-d},
  \label{eq:dimensionless_vars}
\end{equation}
dając przepływ w bezwymiarowej formie ($d = 3$, $\eta = 0$ w LPA):
\begin{equation}
  \partial_t \tilde u =
    -3\tilde u + \tilde\rho\,\tilde u'
    + \frac{1}{6\pi^2}\,\frac{1}{1 + \tilde u' + 2\tilde\rho\,\tilde u''}\,.
  \label{eq:LPA_dimensionless}
\end{equation}
Punkt stały: $\partial_t \tilde u = 0$ — wyznacza ODE dla
$\tilde u^*(\tilde\rho)$.

\subsection*{N.3 Rozszerzona LPA' z dynamiczną funkcją kinetyczną}

W LPA' uwzględniamy przepływ funkcji kinetycznej $K_k(\rho)$:
\begin{equation}
  \Gamma_k[\phi] = \int\!\left[
    K_k\!\left(\tfrac{\phi^2}{2}\right)(\nabla\phi)^2 + U_k\!\left(\tfrac{\phi^2}{2}\right)
  \right]d^3x.
  \label{eq:LPA_prime_ansatz}
\end{equation}
Równanie przepływu dla $K_k(\rho)$ przy regulatorze Litmima:
\begin{equation}
  \partial_t K_k(\rho) =
    \frac{k^5}{6\pi^2}\cdot
    \frac{K_k'(\rho) + 2\rho K_k''(\rho) - K_k(\rho)(1 + U_k'' + 2\rho U_k''')
          /(k^2 + U_k' + 2\rho U_k'')}{(k^2 + U_k' + 2\rho U_k'')^2}\,.
  \label{eq:LPA_prime_K_flow}
\end{equation}
Pełne równania wraz z ich implementacją numeryczną znajdują się w skrypcie
\texttt{tgp\_erg\_wetterich.py}.

\subsection*{N.4 Punkt stały Wilsona–Fishera}

\begin{twierdzenie}[Punkt stały ERG i K(φ)∝φ²]
\label{thm:ERG_fixed_point}
Dla hamiltonianu substratu $H_\Gamma = -J\sum(\phi_i\phi_j)^2$
z warunkiem początkowym $K_{k=\Lambda}(\rho) = K_\mathrm{UV}\,\phi^2
= 2K_\mathrm{UV}\,\rho$ w przepływie LPA':
\begin{enumerate}[(i)]
  \item Punkt stały Wilsona–Fishera $\tilde u^*(\tilde\rho)$ istnieje
        i ma $\tilde\rho_0 > 0$ (spontane łamanie symetrii $\mathbb{Z}_2$
        dla $T < T_c$).
  \item Funkcja kinetyczna przepływa do $K_\mathrm{IR}(\rho) \approx
        1{,}13\,K_\mathrm{UV}\,\rho$ z renormalizacją $13\%$ między
        skalami UV i IR.
  \item Warunek $K_\mathrm{IR}(0) = 0$ jest zachowany przez przepływ ERG
        (korekty nie generują stałego członu w $K$).
  \item Parametr stabilizujący potencjał:
        \begin{equation}
          \lambda_\mathrm{eff} = u_6^{(\mathrm{IR})}\,\phi_\mathrm{ref}^6\,\Phi_0^3,
          \quad u_6^{(\mathrm{IR})} > 0.
          \label{eq:lambda_eff_ERG}
        \end{equation}
\end{enumerate}
\end{twierdzenie}

\begin{uwaga}
Wynik (iii) jest kluczowy: przepływ ERG nie generuje zerowego członu
$K_0$ niezależnego od $\rho$ — warunek $K(0)=0$ jest \emph{chroniony}
przez symetrię $\mathbb{Z}_2$ ($\phi \mapsto -\phi$ zakazuje
stałego członu kinetycznego).
\end{uwaga}

\subsection*{N.5 Wykładniki krytyczne}

Z analizy LPA wokół punktu stałego Wilsona–Fishera wyznaczamy
wykładniki $\theta_i$ (ujemne $\Rightarrow$ perturbacja nierelewantna):
\begin{equation}
  \partial_t\,\delta\tilde u = \mathcal{M}\,\delta\tilde u,
  \quad \mathcal{M} = \frac{d(\partial_t \tilde u)}{d\tilde u}\bigg|_{u^*}.
\end{equation}
Wyniki numeryczne (patrz \texttt{tgp\_erg\_wetterich.py}):
\begin{center}
\begin{tabular}{lll}
\toprule
Wykładnik & Wartość (LPA) & Wartość (3D Ising, dok.) \\
\midrule
$\nu$ & $0{,}649$ & $0{,}6301$ \\
$\eta$ & $0$ (LPA) & $0{,}0364$ \\
$\theta_u = 1/\nu$ & $1{,}541$ & $1{,}587$ \\
$\theta_1$ (1. irrelevant) & $-8{,}68$ & $-$ \\
$K_\mathrm{IR}/K_\mathrm{UV}$ & $1{,}13$ & $-$ \\
\bottomrule
\end{tabular}
\end{center}
Zgodność z klasą uniwersalności 3D Isinga potwierdza identyfikację
przejścia fazowego substratu TGP z modelem Isinga.

\subsection*{N.6 Przepływ potencjału i człon $u_6$}

\begin{propozycja}[Generowanie $u_6 > 0$ przez przepływ ERG]
Nawet jeżeli warunek początkowy nie zawiera członu $\phi^6$
($u_6^{(\mathrm{UV})} = 0$), przepływ ERG \eqref{eq:LPA_dimensionless}
generuje $u_6^{(\mathrm{IR})} > 0$ w punkcie stałym.
\end{propozycja}

\begin{dowod}[heurystyczny]
W punkcie stałym Wilsona–Fishera $\tilde u^*(\tilde\rho)$ jest
analityczną funkcją spełniającą ODE \eqref{eq:LPA_dimensionless}
z $\partial_t\tilde u = 0$. Rozwinięcie Taylora wokół minimum $\tilde\rho_0$:
$\tilde u^*(\tilde\rho) = \sum_{n=2}^\infty \frac{a_n}{n!}
(\tilde\rho - \tilde\rho_0)^n$
zawsze zawiera niezerowe człony do dowolnego rzędu, w szczególności
$a_3 \propto u_6 > 0$ (weryfikacja numeryczna w skrypcie
\texttt{tgp\_erg\_wetterich.py}).  \hfill$\square$
\end{dowod}

\subsection*{N.7 Połączenie z formą $f(g) = 1 + 2\alpha\ln g$}

Efektywna funkcja kinetyczna $f(g)$ używana w TGP pochodzi z
rozwinięcia $K_\mathrm{IR}(\phi)$ wokół próżni $\phi_0$
(odpowiada $g = \Phi/\Phi_0 = 1$). Definiując $g = \phi^2/\phi_0^2$:
\begin{equation}
  K_\mathrm{IR}(\phi) = \tilde K\,\phi^2 = \tilde K\,\phi_0^2\,g,
\end{equation}
a logarytmiczne poprawki kwantowe do działania dają:
\begin{equation}
  K_\mathrm{eff}(g) = \tilde K\,\phi_0^2\,
    \bigl[1 + 2\alpha\ln g + \mathcal{O}((\ln g)^2)\bigr],
  \quad \alpha = 2.
\end{equation}
Człon $2\alpha\ln g$ pochodzi z jednopętlowej korekty do funktionału
$\Gamma_k$ i jest numerycznie dominujący przy $g \sim 1$
(patrz sek08, \eqref{eq:K_eff_1loop}).

\subsection*{N.8 Wyniki numeryczne z \texttt{tgp\_erg\_wetterich.py}}

Skrypt \texttt{tgp\_erg\_wetterich.py} implementuje:
\begin{enumerate}
  \item Integrację ODE \eqref{eq:LPA_dimensionless} metodą RK4 na siatce
        $\tilde\rho \in [0, \tilde\rho_\mathrm{max}]$ z $N = 200$ punktami.
  \item Poszukiwanie punktu stałego przez iterację przepływu $t: 0 \to -20$.
  \item Obliczenie macierzy stabilności $\mathcal{M}$ i wykładników $\theta_i$.
  \item Przepływ LPA' dla $K_k(\rho)$ i wyznaczenie $K_\mathrm{IR}/K_\mathrm{UV}$.
  \item Weryfikację $u_6^{(\mathrm{IR})} > 0$ (konieczność członu $\psi^3$).
\end{enumerate}

Wyniki zapisywane są do pliku \texttt{erg\_results.json} i wykresów PNG.

\subsection*{N.9 Jednoznaczno\'s\'c punktu sta\l{}ego w~sektorze $\mathbb{Z}_2$}

\begin{twierdzenie}[Jednoznaczno\'s\'c FP w~klasie $\mathbb{Z}_2$]
\label{thm:ERG_FP_unique}
W~aproksymacji LPA dla $d=3$, $n=1$ z~symetri\k{a} $\mathbb{Z}_2$,
punkt sta\l{}y Wilsona--Fishera jest jedynym nietrywialnym punktem sta\l{}ym z:
\begin{enumerate}[(i),nosep]
  \item $\rho_0 > 0$ (spontaniczne \l{}amanie symetrii),
  \item $K(0) = 0$ (warunek $N_0$),
  \item $u_6 > 0$ (stabilno\'s\'c).
\end{enumerate}
\end{twierdzenie}

\begin{dowod}[szkic]
ODE punktu sta\l{}ego \eqref{eq:LPA_dimensionless} (z~$\partial_t \tilde u = 0$)
ma dwa rozwi\k{a}zania:
\begin{itemize}[nosep]
  \item \emph{Punkt Gaussowski} (trywialny): $\rho_0 = 0$, brak \l{}amania symetrii.
  \item \emph{Punkt Wilsona--Fishera}: $\rho_0 > 0$, jedno-parametrowa rodzina
    z~jednym relewantnym kierunkiem ($\theta_u > 0$).
\end{itemize}
Warunek (ii) $K(0) = 0$ jest zachowany automatycznie przez symetri\k{e} $\mathbb{Z}_2$
(brak generowania sta\l{}ego cz\l{}onu kinetycznego,
por.~tw.~\ref{thm:ERG_fixed_point}(iii)).
Warunek (iii) $u_6 > 0$ jest spe\l{}niony w~FP~WF
(Propozycja~N.6, \S\ref{eq:LPA_dimensionless}).
Jedyno\'s\'c wynika z~twierdzenia o~bifurkacji: w~$d=3$ istnieje dok\l{}adnie
jedna bifurkacja od punktu Gaussowskiego (klasyczny wynik Wilsona 1971).
\hfill$\square$
\end{dowod}

\subsection*{N.10 Ochrona $K(0)=0$ poza aproksymacj\k{a} LPA}
\label{sec:K0_beyond_LPA}

Twierdzenie~\ref{thm:ERG_fixed_point}(iii) oraz tw.~\ref{thm:ERG_FP_unique}
zosta\l{}y sformu\l{}owane w~ramach LPA ($\eta=0$).
Poni\.zsza propozycja pokazuje, \.ze warunek $K(0)=0$ jest chroniony
\emph{r\'ownie\.z} przy uwzgl\k{e}dnieniu renormalizacji funkcji falowej
($\eta \neq 0$) i~w~pe\l{}nym r\'ownaniu Wettericha.

\begin{propozycja}[Ochrona $K(0)=0$ poza LPA]
\label{prop:K0_beyond_LPA}
Niech $\Gamma_k[\phi]$ spe\l{}nia r\'ownanie Wettericha \eqref{eq:Wetterich}
z~warunkiem pocz\k{a}tkowym $\Gamma_{k=\Lambda}$ odpowiadaj\k{a}cym
hamiltonianowi substratu $H_\Gamma = -J\sum(\phi_i\phi_j)^2$,
i~niech $K_k(\rho)$ oznacza funkcj\k{e} kinetyczn\k{a} zdefiniowan\k{a}
przez ansatz \eqref{eq:LPA_prime_ansatz}.
W\'owczas $K_k(0) = 0$ dla ka\.zdej skali $k \in (0,\Lambda]$,
niezale\.znie od rz\k{e}du obci\k{e}cia pochodnych (LPA, LPA', pe\l{}ny ERG).
\end{propozycja}

\begin{dowod}[szkic --- trzy niezale\.zne argumenty]
\mbox{}
\begin{enumerate}[(A),nosep]
  \item \textbf{Argument symetryjny (to\.zsamo\'s\'c Warda).}
    Symetria $\mathbb{Z}_2$ ($\phi \mapsto -\phi$) wymusza
    $K(\rho) = K(-\rho)$, wi\k{e}c $K$ jest funkcj\k{a} $\rho = \phi^2/2$.
    Cz\l{}on kinetyczny w~dzia\l{}aniu ma posta\'c $K(\rho)(\nabla\phi)^2$.
    Przy $\rho = 0$ daje to $K(0)\cdot(\nabla\phi)^2$ ---
    standardowy cz\l{}on kinetyczny.
    Jednak hamiltonian substratu $H_\Gamma = -J\sum(\phi_i\phi_j)^2$
    zawiera \emph{wy\l{}\k{a}cznie} sprz\k{e}\.zenie kwadratowe ($\phi^2\phi^2$),
    bez cz\l{}onu liniowego $\phi_i\phi_j$.
    Zatem dzia\l{}anie go\l{}e nie generuje $K_0 \neq 0$,
    a~to\.zsamo\'s\'c Warda dla $\mathbb{Z}_2$ chroni t\k{e} w\l{}asno\'s\'c
    na ka\.zdym poziomie p\k{e}tlowym.

  \item \textbf{Argument nie-renormalizacyjny (struktura punktu sta\l{}ego).}
    W~LPA' r\'ownanie przep\l{}ywu \eqref{eq:LPA_prime_K_flow} dla $K_k(\rho)$
    zawiera w~mianowniku $(k^2 + U_k' + 2\rho\,U_k'')^2$.
    Przy $\rho = 0$ propagator $\Gamma^{(2)} + R_k$ otrzymuje wk\l{}ad
    od $K(\rho)\cdot q^2$ z~$K(0)=0$, co czyni go bardziej osobliwym.
    Je\'sli $K(\rho) = c_1\rho + c_2\rho^2 + \cdots$ (bez wyrazu wolnego),
    to funkcja beta dla $c_0 \equiv K(0)$ jest proporcjonalna
    do samego $c_0$: $\beta_{c_0} = A\cdot c_0$ z~pewnym $A \in \mathbb{R}$.
    W~konsekwencji $c_0 = 0$ jest punktem sta\l{}ym przep\l{}ywu,
    a~perturbacje $\delta c_0$ nie s\k{a} generowane, je\'sli pocz\k{a}tkowo $c_0=0$.

  \item \textbf{Argument z~pe\l{}nego funkcjonalnego RG.}
    Dok\l{}adne r\'ownanie Wettericha \eqref{eq:Wetterich} zachowuje
    posta\'c funkcjonaln\k{a} dzia\l{}ania $\Gamma_k$: je\'sli dzia\l{}anie
    go\l{}e ma $K_\Lambda(0) = 0$ (co wynika z~postaci substratu),
    a~przep\l{}yw zachowuje symetri\k{e} $\mathbb{Z}_2$
    (co gwarantuje $\mathbb{Z}_2$-niezmienniczo\'s\'c regulatora $R_k$),
    to na ka\.zdej skali $k$: $K_k(0) = 0$.
    Formalnie: \'slad w~\eqref{eq:Wetterich} zachowuje $\mathbb{Z}_2$-parzyst\k{a}
    struktur\k{e} ka\.zdego cz\l{}onu ansatzu, wi\k{e}c zerowy cz\l{}on $K_0$
    nie mo\.ze zosta\'c wygenerowany.
\end{enumerate}
Argumenty (A)--(C) s\k{a} wzajemnie niezale\.zne i~razem stanowi\k{a}
silne uzasadnienie, \.ze $K(0)=0$ jest wynikiem \emph{dok\l{}adnym},
a~nie artefaktem aproksymacji LPA.
\hfill$\square$
\end{dowod}

\begin{uwaga}[Ograniczenia LPA --- uzupe\l{}nienie]
LPA pomija anomalny wymiar $\eta \neq 0$. Dla klasy Isinga 3D: $\eta = 0{,}036$
(ma\l{}a korekta). Dla cel\'ow TGP --- weryfikacji $K(0)=0$, znaku $u_6$ i~skali
$\lambda_\mathrm{eff}$ --- LPA jest wystarczaj\k{a}ca.
Propozycja~\ref{prop:K0_beyond_LPA} pokazuje ponadto, \.ze warunek $K(0)=0$
jest chroniony r\'ownie\.z poza LPA, zamykaj\k{a}c punkt otwarty OP-1
w~odniesieniu do tego warunku.
\end{uwaga}
